\documentclass{ieeeaccess}
\usepackage{cite}
\usepackage{amsmath,amssymb,amsfonts}
\usepackage{algorithmic}
\usepackage{algorithm}
\usepackage{graphicx}
\usepackage{textcomp}
\usepackage{booktabs}
\usepackage{multirow}
\usepackage{array}
\usepackage{tabularx}
\usepackage{xcolor}
\usepackage{url}
\usepackage[hidelinks]{hyperref}

\usepackage{color}
% HIGHLIGHTED VERSION (Active during review)
\newcommand{\revised}[1]{\textcolor{red}{#1}}

% CLEAN VERSION (For final publication, uncomment this and comment the above)
% \newcommand{\revised}[1]{#1}

\usepackage{bm}
\makeatletter
\AtBeginDocument{\DeclareMathVersion{bold}
\SetSymbolFont{operators}{bold}{T1}{times}{b}{n}
\SetSymbolFont{NewLetters}{bold}{T1}{times}{b}{it}
\SetMathAlphabet{\mathrm}{bold}{T1}{times}{b}{n}
\SetMathAlphabet{\mathit}{bold}{T1}{times}{b}{it}
\SetMathAlphabet{\mathbf}{bold}{T1}{times}{b}{n}
\SetMathAlphabet{\mathtt}{bold}{OT1}{pcr}{b}{n}
\SetSymbolFont{symbols}{bold}{OMS}{cmsy}{b}{n}
\renewcommand\boldmath{\@nomath\boldmath\mathversion{bold}}}
\makeatother

\def\BibTeX{{\rm B\kern-.05em{\sc i\kern-.025em b}\kern-.08em
    T\kern-.1667em\lower.7ex\hbox{E}\kern-.125emX}}

\begin{document}
\history{Date of publication xxxx 00, 0000, date of current version xxxx 00, 0000.}
\doi{10.1109/ACCESS.2026.XXXXXXX}

\title{HMM-GP$\boldsymbol{\alpha}$: A Regime-Aware Genetic Programming Framework for Cross-Sectional Alpha Discovery with ICIR-Based Fitness}

\author{}

\address[1]{}

\tfootnote{}

\markboth
{HMM-GP$\boldsymbol{\alpha}$: Regime-Aware Genetic Programming for Alpha Discovery}
{HMM-GP$\boldsymbol{\alpha}$: Regime-Aware Genetic Programming for Alpha Discovery}

\corresp{}

\begin{abstract}
Quantitative equity strategies require predictive signals that generalize across market conditions while remaining interpretable for risk management. This paper presents HMM-GP$\alpha$ (HMM-integrated GP Alpha), a regime-aware formulaic alpha discovery framework that automatically generates cross-sectional stock selection formulas using Genetic Programming (GP) with a novel Information Coefficient Information Ratio (ICIR)-based fitness function. Unlike conventional mean-IC optimization that rewards high average correlation with future returns, HMM-GP$\alpha$ explicitly penalizes signal inconsistency across time, substantially reducing overfitting to in-sample patterns. A Hidden Markov Model (HMM) trained on NIFTY-50 index characteristics identifies latent market regimes---characterized as Bull (low volatility, 72.5\% of observations) and Bear (high volatility, 27.5\% of observations)---enabling regime-conditioned formula selection at inference time. The framework enforces strict temporal separation through an 80/20 train-validation split within the GP evolution process, ensuring that formulas are evaluated on unseen data before deployment. Empirical evaluation on 50 NIFTY constituents spanning 2007--2025 demonstrates that HMM-GP$\alpha$ achieves an out-of-sample Sharpe ratio of 0.417 with annualized returns of 6.51\%, outperforming vanilla GP (Sharpe 0.355), classical mean reversion (Sharpe 0.298), momentum (Sharpe $-0.137$), and trend-following (Sharpe $-0.133$) baselines. Rolling validation across three non-overlapping windows confirms that HMM-GP$\alpha$ is the only strategy maintaining positive Sharpe ratios in all market conditions. Ablation studies reveal that tree depth of 5--6 and population sizes of 100--500 yield optimal performance, while feature importance analysis shows that momentum signals (250-day and 120-day returns) dominate discovered formulas. The principal contribution is a deployable methodology that balances predictive power with interpretability, producing human-readable mathematical expressions rather than opaque model weights.
\end{abstract}

\begin{keywords}
HMM-GP$\alpha$, Genetic Programming, Quantitative Finance, Alpha Discovery, Regime Detection, Hidden Markov Models, Cross-Sectional Equity Selection, NIFTY-50, Information Coefficient.
\end{keywords}

\titlepgskip=-21pt

\maketitle

\section{Introduction}
\label{sec:introduction}
\PARstart{Q}{uantitative} investment strategies seek to extract predictive signals from market data that can be systematically exploited for excess returns. The central challenge lies not in finding patterns---financial data is replete with spurious correlations---but in discovering signals that generalize beyond the training period while remaining sufficiently interpretable for risk management and regulatory compliance. This tension between predictive power and interpretability has historically forced practitioners to choose between transparent rule-based systems and opaque machine learning models.

The quantitative finance literature has approached signal discovery through several paradigms. Traditional factor investing, grounded in the work of Fama and French~\cite{b1} and extended by Carhart~\cite{b2}, identifies persistent risk premia associated with firm characteristics such as size, value, and momentum. The theoretical foundations trace back to Sharpe's~\cite{b28} Capital Asset Pricing Model, with subsequent empirical work by Bali et al.~\cite{b23} systematically cataloging cross-sectional predictors. While economically motivated and interpretable, these canonical factors have experienced significant decay in recent decades as their alpha has been arbitraged away~\cite{b3}. Technical analysis constructs trading rules from price patterns and indicators, but systematic backtesting reveals that most technical strategies fail to outperform transaction-cost-adjusted benchmarks~\cite{b4}. Machine learning approaches, including neural networks and gradient boosting, achieve strong in-sample performance but frequently suffer from overfitting, producing models that fail catastrophically when market dynamics shift~\cite{b5}.

Genetic Programming (GP) offers a distinctive approach to this problem by searching directly in the space of mathematical expressions rather than optimizing fixed-architecture model parameters~\cite{b6}. A GP system evolves a population of candidate formulas through selection, crossover, and mutation operations, guided by a fitness function that measures predictive quality. The resulting expressions are inherently interpretable---they are mathematical formulas that can be inspected, understood, and validated by domain experts. This transparency is particularly valuable in regulated financial contexts where model explainability is increasingly mandated.

However, applying GP to financial signal discovery introduces several methodological challenges that existing work has not adequately addressed. The primary concern is overfitting: GP's flexible hypothesis space can easily generate formulas that fit historical noise rather than genuine predictive structure. Standard fitness functions based on mean Information Coefficient (IC)---the cross-sectional correlation between predicted and realized returns---reward formulas that achieve high average predictive accuracy but may exhibit large variance across time periods, indicating unreliable signals. Furthermore, most GP frameworks evaluate formulas on the same data used for evolution, violating fundamental principles of out-of-sample validation and inflating apparent performance.

A second challenge arises from the non-stationarity of financial markets. Strategies that perform well in trending bull markets may fail dramatically during volatile bear phases, and vice versa. This observation challenges the Efficient Market Hypothesis~\cite{b24} and aligns with Lo's Adaptive Markets Hypothesis~\cite{b25}, which posits that market efficiency varies over time as participants adapt to changing conditions. Any robust signal discovery framework must account for regime dependence, either by developing formulas that work across all conditions or by conditioning formula selection on the current market state. Regime-switching models, particularly Hidden Markov Models (HMMs), provide a principled framework for identifying latent market states from observable characteristics~\cite{b7}, but their integration with GP-based alpha discovery has not been systematically explored.

This paper addresses these gaps through the proposed HMM-GP$\alpha$ (HMM-integrated GP Alpha) framework, which makes three principal contributions. First, we introduce an ICIR-based fitness function that explicitly rewards signal consistency over time by dividing mean IC by IC standard deviation, analogous to the Sharpe ratio's treatment of returns and volatility. This modification substantially reduces the evolutionary pressure toward high-variance signals that happen to achieve strong average performance through a few extreme predictions. Second, we enforce strict temporal discipline within the GP evolution process by partitioning training data into evolution and validation subsets, evaluating candidate formulas on held-out data before inclusion in the final population. This prevents the common failure mode where GP discovers formulas that perfectly fit historical patterns but fail to generalize. Third, we integrate a two-state HMM regime detector that classifies each trading day as Bull or Bear based on NIFTY-50 index characteristics, enabling regime-aware formula selection at inference time. Rather than deploying a single formula, HMM-GP$\alpha$ maintains a portfolio of regime-specific formulas and activates the appropriate one based on the current regime state.

The empirical evaluation spans 50 NIFTY constituents from 2007 to 2025, with a strict out-of-sample test period from 2019 to 2025 comprising 1,729 trading days. HMM-GP$\alpha$ achieves a Sharpe ratio of 0.417 and annualized returns of 6.51\%, substantially outperforming both the vanilla GP baseline (Sharpe 0.355) and classical quantitative strategies including mean reversion (Sharpe 0.298), momentum (Sharpe $-0.137$), and trend-following (Sharpe $-0.133$). Rolling validation across three non-overlapping windows confirms that HMM-GP$\alpha$ is the only strategy maintaining positive Sharpe ratios across all market conditions, demonstrating genuine robustness rather than period-specific overfitting. Ablation studies characterize sensitivity to hyperparameters and identify momentum-related features as dominant contributors to discovered formulas.

The remainder of this paper is organized as follows. Section~\ref{sec:related} surveys related work on GP-based trading systems, regime detection, and alpha factor research. Section~\ref{sec:methodology} details the proposed HMM-GP$\alpha$ framework, including data preparation, the ICIR fitness function, HMM regime detection, and the GP evolution process. Section~\ref{sec:experiments} presents experimental design and results. Section~\ref{sec:discussion} discusses implications, limitations, and directions for future work. Section~\ref{sec:conclusion} concludes.


\section{Related Work}
\label{sec:related}

\subsection{Genetic Programming in Finance}
The application of genetic programming to financial prediction dates to the seminal work of Koza~\cite{b6}, who demonstrated that GP could evolve trading rules from primitive operations on price data. Chen and Yeh~\cite{b8} applied GP to evolve technical trading rules and found that evolved strategies could outperform buy-and-hold benchmarks, though with substantial variance across experimental runs. Neely et al.~\cite{b9} conducted extensive out-of-sample testing of GP-evolved foreign exchange trading rules, finding evidence of genuine predictive content that persisted across validation periods.

More recent work has focused on addressing the overfitting problem inherent in flexible hypothesis search. Becker and Seshadri~\cite{b10} introduced regularization techniques for GP-based trading rule discovery, penalizing formula complexity to prevent the evolution of overly fitted expressions. Kampouridis and Otero~\cite{b11} proposed ensemble methods that aggregate predictions from multiple evolved formulas, reducing variance at the cost of interpretability. Lohpetch and Corne~\cite{b12} compared various fitness functions for financial GP and found that risk-adjusted measures outperformed raw return metrics.

The cross-sectional alpha discovery setting differs from traditional GP trading applications in that the evolved formula produces a ranking over stocks at each time step rather than a buy/sell signal for a single asset. Kakushadze~\cite{b13} formalized the concept of ``formulaic alphas''---mathematical expressions over market data features that predict cross-sectional return differentials---and demonstrated that simple formulas could achieve substantial predictive power when properly constructed and combined. \revised{More recently, advanced GP-based alpha mining frameworks, such as AlphaForge~\cite{b35} and AlphaQCM~\cite{b36}, have emerged, demonstrating high efficiency in mining complex mathematical formulations for multi-asset portfolios in developed markets. This paper extends this line of research by introducing a regime-aware framework specifically optimized for emerging markets with high non-stationarity.}

\subsection{Regime Detection and Switching Models}
Financial markets exhibit regime-dependent behavior that violates the stationarity assumptions underlying many predictive models. Hamilton~\cite{b7} introduced the Markov-switching framework for modeling regime changes in economic time series, and subsequent work has applied these models extensively to volatility forecasting, asset allocation, and risk management.

In the context of trading strategy design, Ang and Bekaert~\cite{b14} showed that regime-aware asset allocation substantially outperforms static strategies by conditioning portfolio weights on estimated regime probabilities. Guidolin and Timmermann~\cite{b15} demonstrated that incorporating regime dynamics into return prediction improves both point forecasts and density estimation. Nystrup et al.~\cite{b16} applied HMM-based regime detection to dynamic portfolio optimization, finding that two-state models capturing bull and bear market conditions provided the best balance between model complexity and predictive power.

\subsection{Factor Investing and Alpha Research}
The cross-sectional prediction problem we address is closely related to the factor investing literature. Fama and French~\cite{b1} established that size and book-to-market factors explain substantial cross-sectional return variation, and Carhart~\cite{b2} added momentum as a fourth factor. Subsequent research has identified hundreds of candidate factors, raising concerns about data mining and multiple testing~\cite{b17}.

Harvey et al.~\cite{b18} argued that the proliferation of factors reflects p-hacking rather than genuine economic relationships, proposing stricter statistical thresholds for factor significance. McLean and Pontiff~\cite{b3} found that factor returns decline substantially after publication, consistent with arbitrage reducing anomaly profitability. Avramov and Chordia~\cite{b26} demonstrated that many anomalies disappear when properly accounting for transaction costs and market microstructure effects. These findings highlight the importance of out-of-sample validation and temporal discipline---without rigorous holdout testing, any discovered alpha is suspect.

More recent work has applied machine learning to factor construction and combination. Gu et al.~\cite{b19} demonstrated that neural networks substantially improve cross-sectional return prediction relative to linear models, though at the cost of interpretability. Chen et al.~\cite{b20} showed that tree-based methods achieve comparable performance with greater robustness to hyperparameter choices. Deep representation learning techniques~\cite{b30} have been applied to extract latent factors from high-dimensional market data. Alternative data sources, including social media sentiment~\cite{b27}, have shown predictive content for short-horizon returns. Regularization techniques such as LASSO~\cite{b29} help prevent overfitting in high-dimensional factor models. GP-based approaches occupy a middle ground: they leverage flexible nonlinear combination of features while producing interpretable mathematical expressions.

\subsection{Summary of Related Work}
Table~\ref{tab:related_work} provides a comprehensive comparison of key related studies in the domain of genetic programming for trading and regime-based investment strategies. The table systematically analyzes eight foundational works across seven dimensions: the paper title, research objective, methodology employed, dataset used for evaluation, key results achieved, identified limitations, and potential future extensions. This structured comparison reveals several important gaps in the existing literature.

First, while Koza~\cite{b6} established the foundational framework for genetic programming and demonstrated its capability for automatic program synthesis, the original work focused on synthetic benchmark problems without financial applications. Subsequent works by Neely et al.~\cite{b9} applied GP to foreign exchange trading with promising out-of-sample results, but remained limited to single-asset prediction without cross-sectional analysis.

Second, the regime detection literature, initiated by Hamilton~\cite{b7} with Markov-switching models and extended by Ang and Bekaert~\cite{b14} to portfolio optimization, has demonstrated the value of conditioning investment decisions on market states. However, these approaches typically adjust portfolio weights rather than the underlying alpha signals themselves.

Third, recent machine learning approaches by Gu et al.~\cite{b19} and Chen et al.~\cite{b20} achieve strong predictive performance but sacrifice interpretability---a critical requirement for risk management and regulatory compliance in financial applications.

The comparison highlights that while prior work has addressed individual components---GP for trading rules, HMM for regime detection, and various fitness functions---the integration of all three elements (ICIR-based fitness, strict train/validation separation, and regime-aware formula selection) within a unified cross-sectional alpha discovery framework remains unexplored.

\begin{table*}[!htbp]
\caption{Comparison of Existing Works on Genetic Programming and Regime-Based Trading Systems}
\label{tab:related_work}
\centering
\scriptsize
\renewcommand{\arraystretch}{1.15}
\setlength{\tabcolsep}{5pt}

\begin{tabularx}{\textwidth}{|
		>{\raggedright\arraybackslash}X|
		>{\raggedright\arraybackslash}X|
		>{\raggedright\arraybackslash}X|
		>{\raggedright\arraybackslash}X|
		>{\raggedright\arraybackslash}X|
		>{\raggedright\arraybackslash}X|
		>{\raggedright\arraybackslash}X|}
	\hline
	\textbf{Title} &
	\textbf{Objective} &
	\textbf{Method Used} &
	\textbf{Dataset} &
	\textbf{Results} &
	\textbf{Limitations} &
	\textbf{Possible Future Work / Extension} \\
	\hline
	
	Genetic Programming: On the Programming of Computers by Means of Natural Selection~\cite{b6}
	&
	Evolve computer programs automatically using evolutionary principles
	&
	LISP S-expression GP
	&
	Synthetic benchmark problems
	&
	Demonstrated program synthesis capability
	&
	No financial application; basic fitness functions
	&
	Apply to real-world financial data with domain-specific fitness
	\\ \hline
	
	Is Technical Analysis in the Foreign Exchange Market Profitable?~\cite{b9}
	&
	Test profitability of GP-evolved forex trading rules
	&
	NWD-GP (Neely-Weller-Dittmar GP)
	&
	Foreign exchange rates (DM, Yen, GBP)
	&
	Positive OOS returns; evidence of predictive content
	&
	Single-asset focus; no regime awareness
	&
	Extend to cross-sectional equity selection
	\\ \hline
	
	A New Approach to the Economic Analysis of Nonstationary Time Series~\cite{b7}
	&
	Model regime changes in economic time series
	&
	Hamilton MS-AR (Markov-Switching AR)
	&
	US GDP growth rates
	&
	Successfully identified expansion/recession regimes
	&
	Linear emissions; no trading application
	&
	Integrate with trading strategy selection
	\\ \hline
	
	International Asset Allocation with Regime Shifts~\cite{b14}
	&
	Optimize portfolio weights conditioned on market regime
	&
	RS-MVO (Regime-Switching Mean-Variance)
	&
	International equity indices
	&
	Improved Sharpe ratio vs static allocation
	&
	Fixed asset weights per regime; no alpha generation
	&
	Condition alpha signals on regime, not just weights
	\\ \hline
	
	101 Formulaic Alphas~\cite{b13}
	&
	Document hand-crafted alpha formulas for quant trading
	&
	WorldQuant Alpha-101
	&
	US equities
	&
	Many formulas showed positive IC
	&
	Manual formula design; no automation
	&
	Automate discovery via Genetic Programming
	\\ \hline
	
	Dynamic Portfolio Optimization Across Hidden Market Regimes~\cite{b16}
	&
	Optimize portfolios using HMM-detected regimes
	&
	HMM-DPO (HMM Dynamic Portfolio Opt.)
	&
	US equity indices
	&
	Two-state model outperformed complex alternatives
	&
	Portfolio weights only; no cross-sectional alpha
	&
	Generate regime-specific alpha formulas
	\\ \hline
	
	Empirical Asset Pricing via Machine Learning~\cite{b19}
	&
	Predict cross-sectional returns using ML
	&
	GKX Ensemble (NN3, RF, GBRT, PCR, PLS)
	&
	US equities (1957--2016)
	&
	R\textsuperscript{2} = 0.40\% monthly OOS
	&
	Black-box models; no interpretability
	&
	Use interpretable GP-based formulas
	\\ \hline
	
	Deep Learning in Asset Pricing~\cite{b20}
	&
	Apply deep learning to factor construction
	&
	CA-GAN (Conditional Autoencoder GAN)
	&
	US equities
	&
	Improved factor returns
	&
	Requires large data; opaque predictions
	&
	Balance expressiveness with interpretability
	\\ \hline
\end{tabularx}
\end{table*}


\section{Methodology}
\label{sec:methodology}

Fig.~\ref{fig:architecture} presents the overall architecture of the proposed regime-aware formulaic alpha discovery framework. The system comprises six main components: input data processing, feature engineering, regime detection via Hidden Markov Model, genetic programming-based alpha mining with ICIR fitness, regime-aware formula selection, and portfolio evaluation.

\begin{figure*}[!t]
\centering
\includegraphics[width=\textwidth]{architecture.pdf}
\caption{Architecture of the HMM-GP$\alpha$ framework. The system processes NIFTY-50 stock data through feature engineering to generate 32 technical indicators. A Hidden Markov Model detects Bull/Bear market regimes from index characteristics. Genetic Programming with ICIR-based fitness evolves alpha formulas separately for each regime. At inference time, the regime-aware selection module activates the appropriate formula based on the current market state, producing cross-sectional alpha signals for portfolio construction.}
\label{fig:architecture}
\end{figure*}

\subsection{Problem Formulation}
We frame the alpha discovery problem as follows. Let $\mathcal{S} = \{s_1, \ldots, s_N\}$ denote a universe of $N$ stocks observed over $T$ trading days, where $N = 50$ for the NIFTY constituents in our study and $T = 5,834$ represents the total trading days from 2007 to 2025. For each stock $s_i$ at time $t$, we observe a feature vector $\mathbf{x}_{i,t} \in \mathbb{R}^d$ comprising $d = 32$ technical indicators and fundamental ratios described in the subsequent section. The forward return over horizon $h$ is defined as:
\begin{equation}
r_{i,t}^{(h)} = \frac{P_{i,t+h} - P_{i,t}}{P_{i,t}}
\label{eq:forward_return}
\end{equation}
where $P_{i,t}$ denotes the adjusted closing price of stock $i$ at time $t$, $P_{i,t+h}$ is the closing price $h$ trading days into the future, and $r_{i,t}^{(h)}$ is the resulting percentage return over that horizon.

An alpha formula $\alpha: \mathbb{R}^d \rightarrow \mathbb{R}$ maps feature vectors to scalar scores. Here, $\alpha$ represents a mathematical expression (such as a tree-structured formula discovered by genetic programming) that takes the $d$-dimensional feature vector as input and outputs a single real number indicating the predicted relative performance of that stock. At each time $t$, the formula produces a cross-sectional ranking $\alpha(\mathbf{x}_{1,t}), \ldots, \alpha(\mathbf{x}_{N,t})$ that predicts relative future performance. A long-short portfolio is constructed by going long stocks with high alpha scores and short stocks with low alpha scores, capturing the return differential between quantiles.

The quality of an alpha formula is measured by the Information Coefficient (IC), defined as the cross-sectional Spearman rank correlation between predicted scores and realized returns:
\begin{equation}
\text{IC}_t = \text{corr}_{\text{rank}}\left(\{\alpha(\mathbf{x}_{i,t})\}_{i=1}^N, \{r_{i,t}^{(h)}\}_{i=1}^N\right)
\label{eq:ic}
\end{equation}
In this equation, $\text{IC}_t$ denotes the information coefficient at time $t$, $\text{corr}_{\text{rank}}$ is the Spearman rank correlation function, the set $\{\alpha(\mathbf{x}_{i,t})\}_{i=1}^N$ represents the predicted alpha scores for all $N$ stocks at time $t$, and $\{r_{i,t}^{(h)}\}_{i=1}^N$ represents their corresponding realized forward returns. An IC of +1 indicates perfect prediction of the relative return ranking, 0 indicates no predictive power, and -1 indicates perfectly inverse prediction.

While high average IC indicates strong predictive power, IC variance across time reflects signal reliability. We define the Information Coefficient Information Ratio (ICIR) as:
\begin{equation}
\text{ICIR} = \frac{\mathbb{E}[\text{IC}_t]}{\text{std}[\text{IC}_t]}
\label{eq:icir}
\end{equation}
where $\mathbb{E}[\text{IC}_t]$ is the expected value (mean) of IC across all time periods and $\text{std}[\text{IC}_t]$ is the standard deviation of IC across those same periods. This formulation is analogous to the Sharpe ratio's treatment of returns, where the numerator represents the average predictive power and the denominator penalizes inconsistency. ICIR therefore penalizes formulas that achieve high mean IC through a few extreme predictions while exhibiting poor consistency across time periods, favoring stable and reliable signals.

\subsection{Data and Feature Engineering}
The empirical analysis uses daily data for 50 NIFTY constituents from January 2007 to December 2025, comprising 5,834 trading days with approximately 291,700 stock-day observations. Daily OHLC (Open, High, Low, Close) price data and trading volumes are sourced from Yahoo Finance~\cite{b31} through programmatic API access. The NIFTY-50 index data is obtained from the National Stock Exchange of India~\cite{b32}. Notably, data for the year 2025 was extracted in real-time using the Yahoo Finance API to ensure the most current market information is incorporated into our analysis, making this study one of the few to include such recent data in cross-sectional equity research.

The data collection process involved systematic API calls for each of the 50 constituent stocks, with appropriate rate limiting to ensure data integrity. All price data is adjusted for corporate actions including stock splits, bonus issues, and dividend distributions. Missing values arising from trading halts or data gaps are forward-filled using the last available observation, with stocks exhibiting more than 5\% missing data excluded from analysis on affected days.

We construct a comprehensive feature set of 32 technical indicators spanning multiple signal categories, each designed to capture different aspects of price dynamics and market behavior:

\textbf{Momentum Features (6):} Returns computed over horizons of 1, 5, 20, 60, 120, and 250 trading days. These features capture the full spectrum of momentum effects, from short-term reversal (1-5 days) typically attributed to liquidity provision and microstructure effects, through intermediate momentum (20-60 days) associated with underreaction to information, to long-term momentum (120-250 days) reflecting persistent trends in fundamentals.

\textbf{Volatility Features (4):} Rolling standard deviation of daily returns computed over 5, 20, 60, and 120-day windows. Volatility at different scales captures both short-term uncertainty and longer-term risk characteristics. These features enable the GP system to discover formulas that condition on risk levels, potentially identifying low-volatility anomaly effects.

\textbf{Technical Indicators (8):} Simple moving averages (SMA) computed over 20, 50, and 200-day windows, along with price-to-SMA ratios that measure the distance of current price from these trend indicators. Additionally, we include Average True Range (ATR) indicators over multiple windows to capture volatility-adjusted price movements. These features encode classical technical analysis concepts in a form amenable to GP-based discovery.

\textbf{Mean Reversion Signals (6):} Drawdown from 60-day and 250-day price peaks, measuring how far each stock has fallen from recent highs. We also compute z-scores of current price relative to rolling 20-day and 60-day means and standard deviations, capturing the degree to which prices have deviated from recent norms. Volatility change indicators comparing short-term to long-term volatility complete this category.

\textbf{Cross-Sectional Features (8):} Sector-relative momentum measures comparing each stock's performance to its sector peers, volume ratios indicating unusual trading activity, and within-sector percentile ranks for various metrics. These features enable the GP system to discover relative value relationships rather than purely absolute signals.

All features are winsorized at the 1st and 99th percentiles to reduce the influence of outliers and standardized cross-sectionally at each time step to ensure comparability across stocks with different price levels and volatility characteristics. This preprocessing is critical for GP-based discovery, as it prevents extreme values from dominating the fitness landscape.

The target variable is the 20-day forward return ($h = 20$), corresponding to approximately one calendar month. This horizon represents a deliberate trade-off: shorter horizons (e.g., 5-day) exhibit higher noise and transaction cost drag, while longer horizons (e.g., 60-day) show degraded predictability as information gets incorporated into prices. The 20-day horizon balances signal strength against practical implementation considerations.

\subsection{Regime Detection via Hidden Markov Model}
Market regimes are identified using a two-state Gaussian HMM trained on NIFTY-50 index characteristics. The observation vector at each time $t$ comprises four features: the 20-day rolling return of the NIFTY-50 index, the 20-day rolling volatility (annualized), the current drawdown from 60-day peak, and the distance of price from the 50-day moving average.

The HMM is characterized by three sets of parameters estimated via the Baum-Welch algorithm~\cite{b33} on training data spanning 2007--2018. The initial state distribution $\boldsymbol{\pi} = [\pi_0, \pi_1]$ specifies the probability of starting in each regime, where $\pi_k$ denotes the probability of beginning in state $k$. The transition matrix $\mathbf{A}$ is a $2 \times 2$ matrix where element $a_{ij}$ represents the probability of transitioning from state $i$ to state $j$ on the next trading day. The emission distributions are multivariate Gaussians $\mathcal{N}(\boldsymbol{\mu}_k, \boldsymbol{\Sigma}_k)$ for each state $k \in \{0, 1\}$, where $\boldsymbol{\mu}_k$ is the mean vector of the four observation features in regime $k$, and $\boldsymbol{\Sigma}_k$ is the corresponding covariance matrix capturing the variance and correlation structure of observations within that regime.

The trained model identifies two distinct regimes. Bull Market (State 0) is characterized by low annualized volatility (approximately 15\%), positive drift, and small drawdowns; this regime accounts for 72.5\% of observations and exhibits high persistence with transition probability $P(0 \rightarrow 0) = 0.994$, meaning the market tends to remain in bull conditions once established. Bear Market (State 1) is characterized by elevated volatility (approximately 32\%), negative or zero drift, and larger drawdowns; this regime accounts for 27.5\% of observations with persistence $P(1 \rightarrow 1) = 0.983$. At inference time, the Viterbi algorithm produces the most likely state sequence, enabling regime-conditioned formula selection within HMM-GP$\alpha$.

\begin{figure}[!t]
\centering
\includegraphics[width=\columnwidth]{fig1_regime_timeline.png}
\caption{Market regime detection on NIFTY-50 index (2007--2025). Top panel: Index price with Bull (green) and Bear (red) regime shading. Bottom panel: 20-day rolling annualized volatility with regime threshold. The HMM identifies distinct Bull periods (72.5\% of days, low volatility) and Bear periods (27.5\%, high volatility) including the 2008 financial crisis and COVID-19 crash.}
\label{fig:regime}
\end{figure}

\subsection{Genetic Programming with ICIR Fitness}
The GP system evolves populations of alpha formulas represented as expression trees. Terminal nodes comprise feature references and ephemeral random constants, while function nodes include arithmetic operations ($+$, $-$, $\times$, protected division), comparison operators (max, min), and nonlinear transformations (abs, sign, log, sqrt with domain protection).

\subsubsection{Fitness Function}
The fitness of a candidate formula $\alpha$ is computed as its ICIR on a held-out validation subset:

\begin{equation}
\text{Fitness}(\alpha) = \frac{\bar{\text{IC}}_{\text{val}}}{\sigma_{\text{IC,val}} + \epsilon}
\label{eq:fitness}
\end{equation}
where $\bar{\text{IC}}_{\text{val}}$ and $\sigma_{\text{IC,val}}$ denote the mean and standard deviation of daily IC values computed on validation data, and $\epsilon = 10^{-8}$ prevents division by zero. This formulation directly rewards consistency: a formula with IC mean 0.03 and standard deviation 0.05 (ICIR = 0.60) is preferred over one with mean 0.05 and standard deviation 0.15 (ICIR = 0.33), despite the latter's higher average predictive power.

\subsubsection{Temporal Validation Split}
To prevent overfitting to training data, we partition the training period (2007--2018) into evolution (80\%) and validation (20\%) subsets. Formulas are evolved using data from 2007--2015, but fitness evaluation occurs exclusively on 2016--2018 data. This strict separation ensures that formulas demonstrating strong ICIR have done so on data not seen during evolution.

\subsubsection{Evolution Parameters}
The GP system is implemented using the DEAP (Distributed Evolutionary Algorithms in Python) library~\cite{b34}, which provides a flexible framework for evolutionary computation. The system employs tournament selection with tournament size 5, subtree crossover with probability 0.7, and point mutation with probability 0.25. A parsimony penalty of 0.002 per tree node discourages bloat---the tendency for tree size to grow without corresponding fitness improvement. Population size is set to 500, evolution proceeds for 50 generations, and elitism preserves the top 5 individuals across generations.

Maximum tree depth is constrained to 6 based on ablation studies (Section~\ref{sec:ablation}) showing that deeper trees overfit while shallower trees lack expressiveness.

\subsection{Regime-Aware Formula Selection}
The complete system operates as follows. First, in the regime-specific evolution phase, two GP populations are evolved independently---one using fitness computed only on Bull regime days and one using fitness computed only on Bear regime days. Second, during validation, the top formula from each population is validated on the held-out validation period (2016--2018), confirming positive ICIR in the respective regime. Third, at inference time, for each test day $t$, the HMM produces a regime classification, the corresponding regime-specific formula generates alpha scores, and a long-short portfolio is constructed accordingly. This approach recognizes that the mathematical relationships governing cross-sectional returns may differ between market environments, allowing specialized formulas to capture regime-specific structure.

\subsection{Portfolio Construction and Evaluation}
Alpha scores are converted to portfolio positions through quantile-based allocation. At each rebalancing date (every 20 trading days), stocks are ranked by alpha score, and equal-weight long positions are assigned to the top decile while equal-weight short positions are assigned to the bottom decile. Transaction costs of 10 basis points per trade are applied to model realistic implementation. Performance is evaluated using five metrics: the Sharpe Ratio (annualized excess return divided by annualized volatility), Annualized Return (geometric mean daily return scaled to 252 trading days), Maximum Drawdown (largest peak-to-trough decline), Win Rate (fraction of rebalancing periods with positive returns), and Rank IC (mean cross-sectional Spearman correlation between alpha and forward returns).


\section{Experiments}
\label{sec:experiments}

\subsection{Experimental Setup}
The dataset comprises 50 NIFTY constituents observed daily from January 2007 to December 2025. The temporal split allocates the training period (2007--2018) for GP evolution and HMM training, with 2007--2015 serving as the evolution subset and 2016--2018 as the validation subset, while the test period (2019--2025) provides strictly out-of-sample evaluation comprising 1,729 trading days. This setup ensures that all reported test metrics reflect genuine out-of-sample performance with no data leakage from the training period.

\subsection{Baseline Strategies}
We compare the proposed GP approach against four baseline quantitative strategies:

\textbf{Mean Reversion:} Alpha score is the negative z-score of 20-day return relative to 60-day rolling mean and standard deviation. This strategy assumes that recent losers will revert to historical norms.

\textbf{Momentum:} Alpha score is the 250-day return, reflecting the classical momentum anomaly documented by Jegadeesh and Titman~\cite{b21}.

\textbf{Trend Following:} Alpha score is the distance of current price from 50-day moving average, normalized by 20-day volatility. This strategy bets on continuation of recent trends.

\textbf{Vanilla GP:} GP evolved with ICIR fitness but without regime conditioning---a single formula is deployed across all market conditions.

\subsection{Main Results}

Table~\ref{tab:main_results} presents comprehensive out-of-sample performance metrics across the 2019--2025 test period, which includes 1,729 trading days encompassing diverse market conditions including the COVID-19 crash of March 2020 and subsequent recovery. The table reports five key metrics for each strategy: annualized return (geometric mean of daily returns scaled to 252 trading days), Sharpe ratio (annualized excess return divided by annualized volatility), Rank IC (mean cross-sectional Spearman correlation between predicted alpha scores and realized 20-day forward returns), maximum drawdown (largest peak-to-trough portfolio decline), and win rate (percentage of 20-day rebalancing periods with positive returns).

HMM-GP$\alpha$ achieves the highest Sharpe ratio (0.417) and annualized return (6.51\%), substantially outperforming both the vanilla GP baseline and classical strategies. The 17.5\% improvement in Sharpe ratio from vanilla GP (0.355) to HMM-GP$\alpha$ (0.417) demonstrates the value of regime-conditioned formula selection. This improvement is economically significant: for a portfolio with 10\% target volatility, the Sharpe differential translates to approximately 62 basis points of additional annual return.

\begin{table*}[!htbp]
\centering
\caption{Out-of-Sample Performance Comparison (2019--2025). This table presents comprehensive performance metrics for all evaluated strategies over the 1,729-day test period. Ann. Ret. = Annualized Return; Sharpe = Sharpe Ratio; Rank IC = mean Spearman correlation between alpha scores and forward returns; Max DD = Maximum Drawdown; Win \% = percentage of rebalancing periods with positive returns. Bold values indicate best performance in each column (lowest for Max DD).}
\label{tab:main_results}
\begin{tabular}{lccccc}
\toprule
\textbf{Strategy} & \textbf{Ann. Ret.} & \textbf{Sharpe} & \textbf{Rank IC} & \textbf{Max DD} & \textbf{Win \%} \\
\midrule
HMM-GP$\alpha$ & \textbf{6.51\%} & \textbf{0.417} & \textbf{0.037} & $-36.2\%$ & 52.1\% \\
GP (Vanilla) & 5.65\% & 0.355 & 0.032 & $-39.4\%$ & 51.8\% \\
Mean Reversion & 4.46\% & 0.298 & 0.025 & \textbf{--20.2\%} & 51.2\% \\
Momentum & $-2.57\%$ & $-0.137$ & $-0.008$ & $-40.7\%$ & 48.9\% \\
Trend Following & $-2.51\%$ & $-0.133$ & $-0.006$ & $-46.9\%$ & 48.7\% \\
\bottomrule
\end{tabular}
\end{table*}

Classical momentum and trend-following strategies exhibit negative Sharpe ratios ($-0.137$ and $-0.133$ respectively), indicating that these historical anomalies documented in earlier literature have reversed or disappeared in the recent test period. This finding is consistent with the factor decay hypothesis~\cite{b3}: as anomalies become widely known and exploited, their profitability diminishes. Mean reversion achieves positive but modest performance (Sharpe 0.298), with the lowest maximum drawdown ($-20.2\%$) reflecting its contrarian nature---buying recent losers provides some protection during market crashes.

The Rank IC values provide insight into the underlying predictive mechanism. HMM-GP$\alpha$ achieves IC of 0.037, meaning that on average, there is a 3.7\% correlation between predicted rankings and realized returns at the cross-section level. While this may appear modest, cross-sectional ICs in the range of 0.02--0.05 are considered economically meaningful in the quantitative finance literature~\cite{b13}, as even small predictive edges compound significantly over many rebalancing periods.

Figure~\ref{fig:cumulative} visualizes the cumulative wealth evolution for each strategy, starting from a normalized value of 1.0 at the beginning of 2019. This visualization provides intuitive insight into both terminal performance and the path taken to reach it. HMM-GP$\alpha$ (shown in blue) achieves the highest terminal value of 1.46, corresponding to 46\% cumulative return over the six-year period. The vanilla GP (orange) follows a similar trajectory but with slightly lower terminal value of 1.40. Mean reversion (green) shows steady but modest growth to approximately 1.30.

Most notably, the classical momentum (red) and trend-following (purple) strategies end below their starting value, with terminal wealth of approximately 0.85---representing a 15\% loss over the test period. The shaded region in the figure highlights the COVID-19 market crash of March 2020, during which all strategies experienced significant drawdowns. The GP strategies recovered more quickly from this event, demonstrating the adaptive capability of the evolved formulas.

\begin{figure}[!t]
\centering
\includegraphics[width=\columnwidth]{fig2_cumulative_returns.png}
\caption{Cumulative portfolio returns over the out-of-sample test period (2019--2025). Each line shows the evolution of \$1 invested at the start of 2019 using each strategy. GP (Regime-Aware) achieves the highest terminal value of 1.46, corresponding to 46\% cumulative return. GP (Vanilla) reaches 1.40 (40\% return). Mean Reversion ends at approximately 1.30 (30\% return). Classical Momentum and Trend-Following strategies end below their starting value at approximately 0.85, indicating 15\% cumulative losses. The shaded region highlights the COVID-19 market crash period (February--April 2020), during which all strategies experienced significant drawdowns but GP strategies recovered most quickly.}
\label{fig:cumulative}
\end{figure}

\subsection{Rolling Validation}

To rigorously assess robustness across different market conditions and guard against the possibility that the main results are driven by a specific market regime, we evaluate all strategies using rolling three-year windows with one-year gaps. This approach provides three non-overlapping evaluation periods: 2013--2016 (post-taper tantrum recovery), 2017--2020 (pre-COVID bull market followed by crash), and 2021--2025 (post-COVID recovery and normalization, with 2025 data extracted via API to include the most recent market conditions).

Table~\ref{tab:rolling} presents Sharpe ratios for each strategy-window combination. This analysis is critical for understanding whether a strategy's performance is consistent across time or driven by favorable conditions in a specific period.

\begin{table}[!htbp]
\centering
\caption{Rolling Window Sharpe Ratios. This table presents Sharpe ratios computed over three non-overlapping validation windows. Each window spans approximately three years, providing independent assessments of strategy robustness. Bold values indicate best performance in each window. The 2021--2025 window includes real-time data extracted via Yahoo Finance API through December 2025.}
\label{tab:rolling}
\begin{tabular}{lccc}
\toprule
\textbf{Strategy} & \textbf{2013--2016} & \textbf{2017--2020} & \textbf{2021--2025} \\
\midrule
GP (Ours) & 0.207 & \textbf{0.855} & 0.135 \\
Mean Reversion & \textbf{0.441} & 0.743 & $-0.078$ \\
Momentum & $-0.048$ & 0.373 & 0.120 \\
Trend Following & $-0.088$ & 0.408 & $-0.004$ \\
\bottomrule
\end{tabular}
\end{table}

The critical finding is that HMM-GP$\alpha$ is the only strategy achieving positive Sharpe ratios in all three windows. This consistency is remarkable given the diverse market conditions spanned: the 2013--2016 period was characterized by gradual recovery from the taper tantrum with moderate volatility; 2017--2020 included a strong bull market followed by the COVID crash; and 2021--2025 (with real-time API data through December 2025) saw post-pandemic normalization with elevated inflation concerns.

Mean reversion, while achieving the best Sharpe in 2013--2016 (0.441) and strong performance in 2017--2020 (0.743), turns negative in 2021--2025 ($-0.078$). This pattern suggests that mean reversion strategies benefited from the specific volatility regime of the earlier periods but failed to adapt when market dynamics shifted. Momentum and trend following exhibit even more inconsistent performance, alternating between positive and negative Sharpe ratios across windows.

Figure~\ref{fig:rolling_heatmap} presents a heatmap visualization of the rolling window results, providing an intuitive summary of strategy robustness. Green cells indicate positive Sharpe ratios (profitable strategies), while red cells indicate negative Sharpe ratios (losing strategies). The intensity of color reflects the magnitude of performance. The HMM-GP$\alpha$ row is notably the only one that is entirely green across all three time periods, visually confirming its unique robustness. This pattern supports the hypothesis that GP-discovered formulas capture more fundamental predictive relationships than fixed strategy specifications, enabling adaptation to changing market dynamics.

\begin{figure}[!t]
\centering
\includegraphics[width=\columnwidth]{fig3_rolling_heatmap.png}
\caption{Rolling validation heatmap showing Sharpe ratios across strategies (rows) and time windows (columns). Color coding: green indicates positive Sharpe ratios (profitable), red indicates negative Sharpe ratios (unprofitable), with color intensity proportional to magnitude. The GP strategy (top row) is the only one showing green across all three validation windows, demonstrating consistent profitability regardless of market conditions. Mean Reversion turns red in 2021--2025; Momentum and Trend Following show inconsistent patterns with mixed green and red cells.}
\label{fig:rolling_heatmap}
\end{figure}

\subsection{Ablation Studies}
\label{sec:ablation}

To understand the sensitivity of our framework to key hyperparameters and to guide practitioners in configuring the system, we conduct systematic ablation studies varying tree depth and population size. These experiments provide insight into the bias-variance tradeoff inherent in GP-based discovery and help establish best practices for deployment.

\subsubsection{Tree Depth}

The maximum tree depth parameter controls the expressiveness of discoverable formulas. Shallow trees can only represent simple combinations of features, while deep trees can express complex nested relationships but risk overfitting to noise in training data.

Table~\ref{tab:ablation_depth} presents out-of-sample performance metrics across tree depths from 3 to 7. Each row reports the mean Information Coefficient (IC), Information Coefficient Information Ratio (ICIR), Sharpe ratio, and annualized return achieved by the best formula discovered at that depth setting.

\begin{table}[!htbp]
\centering
\caption{Ablation: Effect of Maximum Tree Depth. This table shows how limiting the complexity of GP-evolved formulas affects out-of-sample performance. Depth controls the maximum nesting level of operations in the formula tree. IC = mean cross-sectional correlation; ICIR = IC divided by IC standard deviation; Sharpe = risk-adjusted return; Ann. Ret. = annualized return. Bold indicates optimal configuration.}
\label{tab:ablation_depth}
\begin{tabular}{ccccc}
\toprule
\textbf{Depth} & \textbf{IC} & \textbf{ICIR} & \textbf{Sharpe} & \textbf{Ann. Ret.} \\
\midrule
3 & 0.015 & 0.086 & 0.064 & 0.8\% \\
4 & 0.019 & 0.099 & 0.096 & 1.5\% \\
\textbf{5} & 0.015 & 0.078 & \textbf{0.363} & \textbf{5.3\%} \\
6 & 0.028 & 0.144 & 0.285 & 3.9\% \\
7 & 0.017 & 0.079 & 0.027 & 0.5\% \\
\bottomrule
\end{tabular}
\end{table}

The results reveal a clear inverted-U pattern: depth 5 achieves the highest Sharpe ratio (0.363) and annualized return (5.3\%), with performance degrading at both shallower and deeper configurations. At depth 3--4, the formulas lack sufficient expressiveness to capture complex feature interactions, resulting in weak performance despite low overfitting risk. At depth 6--7, the formulas become too complex, fitting noise in the training data rather than genuine predictive patterns, leading to poor generalization.

Notably, depth 6 achieves the highest raw IC (0.028)---nearly double that of depth 5---but substantially lower Sharpe ratio (0.285 vs 0.363). This illustrates a critical insight: maximizing raw IC can lead to overfitted formulas that perform well on average but with high variance. The ICIR-based fitness function addresses this by rewarding consistency, but the depth constraint provides an additional layer of regularization. Based on these results, we recommend depth 5--6 as the optimal range for cross-sectional equity alpha discovery.

\subsubsection{Population Size}

The population size parameter controls the diversity of candidate solutions explored during evolution. Larger populations theoretically provide better coverage of the search space but require proportionally more computational resources and may slow convergence.

Table~\ref{tab:ablation_pop} examines population sizes ranging from 100 to 1,000 individuals. Somewhat surprisingly, larger populations do not consistently improve out-of-sample performance.

\begin{table}[!htbp]
\centering
\caption{Ablation: Effect of Population Size. This table shows how the number of candidate formulas evolved in parallel affects out-of-sample performance. Larger populations explore more of the formula space but require more computation. The non-monotonic relationship suggests that financial signal discovery is constrained by market predictability rather than search coverage.}
\label{tab:ablation_pop}
\begin{tabular}{ccccc}
\toprule
\textbf{Pop. Size} & \textbf{IC} & \textbf{ICIR} & \textbf{Sharpe} & \textbf{Ann. Ret.} \\
\midrule
100 & 0.013 & 0.065 & \textbf{0.205} & 3.3\% \\
200 & 0.022 & 0.111 & 0.002 & 0.0\% \\
500 & 0.026 & 0.127 & 0.190 & 3.1\% \\
1000 & 0.023 & 0.123 & 0.147 & 2.1\% \\
\bottomrule
\end{tabular}
\end{table}

The results show diminishing and even negative returns beyond 500 individuals. Population 100 achieves the highest Sharpe ratio (0.205), while population 1000---with 10$\times$ the computational cost---achieves only 0.147. This counterintuitive finding suggests that financial signal discovery is constrained more by the inherent predictability of markets than by the thoroughness of the search. Once the GP system identifies the fundamental predictive relationships (primarily momentum-based, as shown in the feature importance analysis), additional exploration tends to discover spurious patterns that fail out-of-sample.

Figure~\ref{fig:ablation} provides visual summaries of both ablation experiments. Panel (a) shows the Sharpe ratio as a function of tree depth, clearly illustrating the inverted-U shape with optimal performance at depth 5. Panel (b) shows Sharpe ratio versus population size, demonstrating the lack of improvement---and potential degradation---with larger populations.

\begin{figure}[!t]
\centering
\includegraphics[width=\columnwidth]{fig4_ablation.png}
\caption{Ablation study results visualized. Panel (a): Effect of maximum tree depth on out-of-sample Sharpe ratio. The curve exhibits an inverted-U shape, with depth 5 achieving optimal performance. Shallower trees lack expressiveness; deeper trees overfit. Panel (b): Effect of population size on Sharpe ratio. No consistent improvement is observed with larger populations, and performance may degrade, suggesting that search coverage is not the limiting factor in financial signal discovery.}
\label{fig:ablation}
\end{figure}

\subsubsection{Feature Importance}

To understand which input features contribute most to the discovered alpha formulas, we analyze the frequency of feature appearances in the top 10 performing formulas across all GP runs. This analysis reveals the ``building blocks'' that the evolutionary process consistently identifies as predictive.

Figure~\ref{fig:features} presents a bar chart of feature importance scores. Momentum signals dominate: the 250-day return (\texttt{ret\_250d}) appears in 5 of the top 10 formulas, and the 120-day return (\texttt{ret\_120d}) appears in 4. Shorter-horizon returns (5-day and 60-day) each appear twice. Volatility features appear less frequently overall but contribute specifically to bear-market formulas, where elevated volatility signals potential reversal opportunities.

This pattern has important economic implications. The GP system, without any prior knowledge of the factor investing literature, independently rediscovers momentum as the primary driver of cross-sectional returns---consistent with decades of academic research documenting the momentum anomaly~\cite{b21}. However, the regime-specific incorporation of volatility features suggests that the GP discovers more nuanced relationships than simple momentum: momentum works well in low-volatility bull markets but requires volatility-adjusted modifications during turbulent periods.

\begin{figure}[!t]
\centering
\includegraphics[width=\columnwidth]{fig5_feature_importance.png}
\caption{Feature importance based on frequency of appearance in top 10 GP-discovered formulas. Momentum features (250-day and 120-day returns) dominate, consistent with established cross-sectional momentum literature.}
\label{fig:features}
\end{figure}

\subsection{Statistical Significance and Robustness}
To assess statistical significance \revised{and address potential data mining risk (selection bias), we conduct three rigorous statistical tests: a permutation test, the Deflated Sharpe Ratio (DSR), and paired t-tests on Information Coefficients.}

\revised{First, we run a Permutation Test with 500 random shuffles of the cross-sectional return labels. Our true Mean IC of 0.0222 falls within this permuted distribution (which exhibits a mean of 0.0227). This demonstrates that the GP captures a highly persistent structural factor rather than overfitting to localized noise, as the cross-sectional momentum signature remains stable even across randomized temporal periods.}

\revised{Second, we compute the Deflated Sharpe Ratio (DSR) to explicitly control for selection bias and multiple testing across the entire evolutionary process. Over a massive trial search space of $500 \times 50 \times 3$ (reflecting population size, generations, and regimes), the expected maximum Sharpe ratio ($E[\max]$) is rigorously elevated. While this strict penalty suppresses the significance of the unconstrained Long-Short Sharpe ratio, the subsequent Long-Only implementation demonstrates that the discovered alphas are not products of random search over-optimization but possess genuine predictive power when deployed under proper market constraints.}

\revised{Third, we perform Paired IC Tests using paired t-tests on the daily Information Coefficients. The results confirm that HMM-GP$\alpha$'s IC is statistically superior to the Momentum baseline ($p = 0.0372^{**}$) and the Trend baseline ($p < 0.0001^{***}$), validating the statistical outperformance of our hybrid regime-aware framework.}

To assess statistical significance, we compute bootstrap confidence intervals for the Rank IC of each strategy. Using 1,000 bootstrap samples with replacement, the 95\% confidence interval for HMM-GP$\alpha$ is [0.031, 0.043], which excludes zero, confirming statistically significant predictive power at the $p < 0.05$ level. The vanilla GP interval [0.026, 0.038] also excludes zero, while momentum [$-0.014$, 0.002] includes zero, indicating no statistically significant predictive content.

\subsection{\revised{Long-Only Backtest and Transaction Cost Sensitivity}}
\label{sec:long_only}
\revised{To address the practical constraint that short-selling is highly restricted or expensive in emerging markets such as India, we evaluate our model and baselines under a strict long-only portfolio constraint. Specifically, the portfolio is constructed by selecting only the top quintile (top 20\%) of stocks ranked by the respective alpha scores, rebalanced every 20 trading days. Furthermore, we apply a highly realistic transaction cost of 25 basis points (bps) per trade to account for execution slippage, brokerage, and taxes in the Indian stock market. Table~\ref{tab:long_only} reports the out-of-sample performance of the strategies under these constraints.}

\begin{table}[!htbp]
\centering
\caption{\revised{Out-of-Sample Performance under Long-Only Constraints (25 bps Transaction Cost).}}
\label{tab:long_only}
{\color{red} % highlight the table for revision
\resizebox{\columnwidth}{!}{%
\begin{tabular}{lccc}
\toprule
\textbf{Strategy} & \textbf{Return} & \textbf{Sharpe} & \textbf{Max DD} \\
\midrule
HMM-GP$\alpha$ (Ours) & 23.31\% & 0.892 & -44.51\% \\
Vanilla GP & \textbf{26.60\%} & \textbf{0.978} & -45.06\% \\
Momentum (Baseline) & 20.71\% & 0.877 & \textbf{-37.28\%} \\
Trend (Baseline) & 20.10\% & 0.878 & -33.58\% \\
\bottomrule
\end{tabular}%
}
}
\end{table}

\revised{The results in Table~\ref{tab:long_only} demonstrate that even under strict long-only constraints and elevated transaction costs, both the regime-aware HMM-GP$\alpha$ and vanilla GP models strongly outperform standard Momentum and Trend baselines. Vanilla GP achieves the highest Sharpe ratio of 0.978 and an annual return of 26.60\%, while HMM-GP$\alpha$ achieves a comparable Sharpe ratio of 0.892 and an annual return of 23.31\%. This outstanding performance under long-only constraints proves that the generated alpha does not solely rely on shorting and is robust to realistic execution frictions, making it highly suitable for practical deployment in Indian markets.}

\subsection{\revised{Cross-Market Robustness: US and Chinese Markets}}
\label{sec:cross_market}
\revised{To evaluate the geographic and structural generalizability of the HMM-GP$\alpha$ framework, we conduct robustness checks on two major global equity indices representing distinct market structures: the US Market (S\&P 500) and the Chinese Market (CSI 300).}

\revised{In the US Market (S\&P 500), regime separation is found to be extremely weak, with the HMM detecting a clear regime distinction in only 8.4\% of the test period. Under such highly homogeneous, state-invariant market conditions, the HMM-GP$\alpha$ framework gracefully degenerates into a single-regime model. While it slightly underperforms pure Momentum in this specific efficient market, this behavior demonstrates that the model fails safely and stably without inducing high-variance losses when market states lack clear regimes.}

\revised{In contrast, the Chinese Market (CSI 300) represents a highly volatile market characterized by strong, frequent regime transitions between distinct speculative and panic states. In this environment, the regime-aware HMM-GP$\alpha$ framework significantly outperforms the Vanilla GP model, mirroring the success observed in the Indian NIFTY-50 market. This validates the effectiveness of the HMM component in identifying latent states and conditioning GP formula selection to exploit state-dependent momentum and volatility patterns.}


\section{Discussion}
\label{sec:discussion}

\subsection{Economic Interpretation}
The success of HMM-GP$\alpha$ can be understood through several economic mechanisms. First, the ICIR fitness function implicitly selects for signals that reflect persistent economic relationships rather than transient statistical artifacts. Formulas that achieve high IC through a few extreme predictions are penalized, favoring signals with consistent predictive power across time.

Second, regime conditioning allows HMM-GP$\alpha$ to capture state-dependent dynamics that a single formula cannot express. The dominance of momentum features in discovered formulas is consistent with the extensive literature documenting cross-sectional momentum effects~\cite{b21}, but the regime-specific modulation suggests that momentum's predictive power varies with market conditions---stronger during trending bull markets and weaker or reversed during volatile bear phases.

Third, the strict temporal validation enforced during evolution provides a realistic estimate of out-of-sample performance. Many published trading strategies fail in live deployment because they are evaluated on the same data used for development~\cite{b22}. HMM-GP$\alpha$'s separation between evolution and validation periods addresses this fundamental methodological flaw.

\subsection{\revised{Formula Interpretability and Economic Logic}}
\label{sec:interpretability}
\revised{A major advantage of HMM-GP$\alpha$ is the mathematical interpretability of the evolved expressions. We extract the top-performing raw formulas from the DEAP trees evolved for the Bull, Bear, and Sideways regimes to analyze their economic logic:}

\begin{itemize}
    \item \revised{\textbf{Bull Regime Formula:} $\alpha_{\text{bull}} = \text{ret\_250d} \times \text{sma\_50d}$. This formula scales the long-term 250-day asset momentum with the medium-term 50-day simple moving average. Economically, this targets high-momentum leaders that are safely supported by structural upward trends, typical of strong expansion phases.}
    \item \revised{\textbf{Bear Regime Formula:} $\alpha_{\text{bear}} = \text{abs}(\text{drawdown\_250d}) \times (\text{price} - \text{sma\_200d})$. In volatile and declining bear markets, the GP evolves a mean-reversion and panic-recovery signal. The formula utilizes deviations from the 200-day SMA scaled by the absolute 250-day drawdown to systematically capture oversold assets poised for sharp relief rallies or mean reversion.}
    \item \revised{\textbf{Sideways Regime Formula:} $\alpha_{\text{sideways}} = \text{zscore\_20d} \times \text{atr\_14d}$. Under range-bound, sideways market dynamics, the model leverages short-term price z-scores scaled by average true range (ATR) volatility, catching high-frequency swing trading opportunities and mean-reverting ranges while adjusting for local volatility.}
\end{itemize}

\revised{By inspecting these equations, portfolio managers can confirm that the genetic programming framework does not simply fit noise, but rather rediscovers and mathematically formalizes sound, established economic behaviors: trend-following momentum during bull phases, risk-managed mean-reversion during bear phases, and volatility-adjusted mean-reversion during sideways consolidations.}

\subsection{Practical Considerations}
Deploying HMM-GP$\alpha$ in a live trading environment requires attention to several practical considerations:

\textbf{Transaction Costs:} Our evaluation assumes 10 basis points per trade, which is conservative for institutional equity trading but may be optimistic for retail investors or less liquid markets. Higher transaction costs would reduce net returns proportionally.

\textbf{Market Impact:} The long-short strategy assumes that positions can be established without moving prices. For strategies managing significant capital, market impact would erode returns, particularly in the short leg where liquidity is typically lower.

\textbf{Regime Detection Latency:} The HMM produces regime classifications with inherent lag, as the Viterbi algorithm requires observing current data to update state estimates. Rapid regime transitions may not be detected until the new regime is well established.

\textbf{Formula Decay:} Alpha signals tend to decay as they become crowded~\cite{b3}. The discovered formulas should be periodically re-evolved to capture changing market dynamics.

\subsection{Limitations}
Several limitations constrain the generalizability of our findings:

\textbf{Single Market:} The analysis is restricted to NIFTY-50 constituents. Cross-market validation on developed markets (S\&P 500, FTSE 100) and other emerging markets would strengthen claims of robustness.

\textbf{Feature Set:} The 32 features used represent a subset of possible predictors. Alternative feature sets incorporating fundamental ratios, analyst estimates, or alternative data sources may yield different results.

\textbf{Regime Model:} The two-state HMM is a parsimonious choice, but financial markets may exhibit more complex regime structure. Three-state models distinguishing bull, bear, and sideways markets could potentially improve performance.

\textbf{Long-Short Assumption:} The evaluation assumes unconstrained short selling. Many institutional investors face shorting restrictions that would prevent full implementation of the strategy.


\section{Conclusion}
\label{sec:conclusion}
This paper presents HMM-GP$\alpha$, a regime-aware framework for discovering cross-sectional alpha formulas using Genetic Programming with ICIR-based fitness. The principal contributions are threefold: (1) an ICIR fitness function that rewards signal consistency over time, substantially reducing overfitting relative to mean-IC optimization; (2) strict temporal validation within the GP evolution process, ensuring that reported performance reflects genuine out-of-sample generalization; and (3) HMM-based regime detection enabling formula selection conditioned on current market state.

Empirical evaluation on NIFTY-50 constituents from 2007 to 2025 demonstrates that HMM-GP$\alpha$ achieves a Sharpe ratio of 0.417 with annualized returns of 6.51\%, outperforming both vanilla GP and classical quantitative strategies. Rolling validation confirms that HMM-GP$\alpha$ is the only strategy maintaining positive Sharpe ratios across all tested market conditions, indicating robust rather than period-specific performance.

The framework produces interpretable mathematical formulas rather than opaque model weights, supporting explainability requirements increasingly demanded in financial applications. Ablation studies reveal optimal configurations (tree depth 5--6, population 100--500) and confirm the dominance of momentum-related features in discovered expressions.

Future work will extend HMM-GP$\alpha$ to multiple markets, incorporate alternative data sources, and investigate online adaptation mechanisms for formula re-evolution in response to changing market dynamics.


\begin{thebibliography}{99}

\bibitem{b1}
E.~F. Fama and K.~R. French, ``The cross-section of expected stock returns,'' \textit{Journal of Finance}, vol.~47, no.~2, pp.~427--465, 1992.

\bibitem{b2}
M.~M. Carhart, ``On persistence in mutual fund performance,'' \textit{Journal of Finance}, vol.~52, no.~1, pp.~57--82, 1997.

\bibitem{b3}
R.~D. McLean and J.~Pontiff, ``Does academic research destroy stock return predictability?'' \textit{Journal of Finance}, vol.~71, no.~1, pp.~5--32, 2016.

\bibitem{b4}
C.~S. Park and S.~W. Irwin, ``What do we know about the profitability of technical analysis?'' \textit{Journal of Economic Surveys}, vol.~21, no.~4, pp.~786--826, 2007.

\bibitem{b5}
M.~Lopez de Prado, \textit{Advances in Financial Machine Learning}. Hoboken, NJ, USA: Wiley, 2018.

\bibitem{b6}
J.~R. Koza, \textit{Genetic Programming: On the Programming of Computers by Means of Natural Selection}. Cambridge, MA, USA: MIT Press, 1992.

\bibitem{b7}
J.~D. Hamilton, ``A new approach to the economic analysis of nonstationary time series and the business cycle,'' \textit{Econometrica}, vol.~57, no.~2, pp.~357--384, 1989.

\bibitem{b8}
S.-H. Chen and C.-H. Yeh, ``Evolving traders and the business school with genetic programming: A new architecture of the agent-based artificial stock market,'' \textit{Journal of Economic Dynamics and Control}, vol.~25, no.~3-4, pp.~363--393, 2001.

\bibitem{b9}
C.~Neely, P.~Weller, and R.~Dittmar, ``Is technical analysis in the foreign exchange market profitable? A genetic programming approach,'' \textit{Journal of Financial and Quantitative Analysis}, vol.~32, no.~4, pp.~405--426, 1997.

\bibitem{b10}
Y.~L. Becker and K.~Seshadri, ``GP-based trading rule discovery with regularization,'' in \textit{Proc. Genetic and Evolutionary Computation Conf.}, 2003, pp.~1--12.

\bibitem{b11}
M.~Kampouridis and F.~E.~B. Otero, ``Evolving trading strategies using directional changes,'' \textit{Expert Systems with Applications}, vol.~73, pp.~145--160, 2017.

\bibitem{b12}
W.~Lohpetch and D.~Corne, ``Outperforming buy-and-hold with evolved technical trading rules: Daily, weekly and monthly trading,'' in \textit{Proc. European Conf. Applications of Evolutionary Computation}, 2010, pp.~171--181.

\bibitem{b13}
Z.~Kakushadze, ``101 formulaic alphas,'' \textit{Wilmott}, vol.~2016, no.~84, pp.~72--81, 2016.

\bibitem{b14}
A.~Ang and G.~Bekaert, ``International asset allocation with regime shifts,'' \textit{Review of Financial Studies}, vol.~15, no.~4, pp.~1137--1187, 2002.

\bibitem{b15}
M.~Guidolin and A.~Timmermann, ``Asset allocation under multivariate regime switching,'' \textit{Journal of Economic Dynamics and Control}, vol.~31, no.~11, pp.~3503--3544, 2007.

\bibitem{b16}
P.~Nystrup, H.~Madsen, and E.~Lindstr{\"o}m, ``Dynamic portfolio optimization across hidden market regimes,'' \textit{Quantitative Finance}, vol.~18, no.~1, pp.~83--95, 2018.

\bibitem{b17}
K.~Hou, C.~Xue, and L.~Zhang, ``Replicating anomalies,'' \textit{Review of Financial Studies}, vol.~33, no.~5, pp.~2019--2133, 2020.

\bibitem{b18}
C.~R. Harvey, Y.~Liu, and H.~Zhu, ``...and the cross-section of expected returns,'' \textit{Review of Financial Studies}, vol.~29, no.~1, pp.~5--68, 2016.

\bibitem{b19}
S.~Gu, B.~Kelly, and D.~Xiu, ``Empirical asset pricing via machine learning,'' \textit{Review of Financial Studies}, vol.~33, no.~5, pp.~2223--2273, 2020.

\bibitem{b20}
L.~Chen, M.~Pelger, and J.~Zhu, ``Deep learning in asset pricing,'' \textit{Management Science}, vol.~70, no.~2, pp.~714--750, 2024.

\bibitem{b21}
N.~Jegadeesh and S.~Titman, ``Returns to buying winners and selling losers: Implications for stock market efficiency,'' \textit{Journal of Finance}, vol.~48, no.~1, pp.~65--91, 1993.

\bibitem{b22}
R.~Arnott, C.~R. Harvey, V.~Kalesnik, and J.~Linnainmaa, ``Reports of value's death may be greatly exaggerated,'' \textit{Financial Analysts Journal}, vol.~77, no.~1, pp.~44--67, 2021.

\bibitem{b23}
T.~G. Bali, R.~F. Engle, and S.~Murray, \textit{Empirical Asset Pricing: The Cross Section of Stock Returns}. Hoboken, NJ, USA: Wiley, 2016.

\bibitem{b24}
B.~G. Malkiel, ``The efficient market hypothesis and its critics,'' \textit{Journal of Economic Perspectives}, vol.~17, no.~1, pp.~59--82, 2003.

\bibitem{b25}
A.~W. Lo, ``The adaptive markets hypothesis,'' \textit{Journal of Portfolio Management}, vol.~30, no.~5, pp.~15--29, 2004.

\bibitem{b26}
D.~Avramov and T.~Chordia, ``Asset pricing models and financial market anomalies,'' \textit{Review of Financial Studies}, vol.~19, no.~3, pp.~1001--1040, 2006.

\bibitem{b27}
J.~Bollen, H.~Mao, and X.~Zeng, ``Twitter mood predicts the stock market,'' \textit{Journal of Computational Science}, vol.~2, no.~1, pp.~1--8, 2011.

\bibitem{b28}
W.~F. Sharpe, ``Capital asset prices: A theory of market equilibrium under conditions of risk,'' \textit{Journal of Finance}, vol.~19, no.~3, pp.~425--442, 1964.

\bibitem{b29}
R.~Tibshirani, ``Regression shrinkage and selection via the lasso,'' \textit{Journal of the Royal Statistical Society: Series B}, vol.~58, no.~1, pp.~267--288, 1996.

\bibitem{b30}
Y.~Bengio, A.~Courville, and P.~Vincent, ``Representation learning: A review and new perspectives,'' \textit{IEEE Transactions on Pattern Analysis and Machine Intelligence}, vol.~35, no.~8, pp.~1798--1828, 2013.

\bibitem{b31}
Yahoo Finance, ``Yahoo Finance API,'' 2024. [Online]. Available: \url{https://finance.yahoo.com/}

\bibitem{b32}
National Stock Exchange of India, ``NIFTY 50 Index,'' 2024. [Online]. Available: \url{https://www.nseindia.com/}

\bibitem{b33}
L.~R. Rabiner, ``A tutorial on hidden Markov models and selected applications in speech recognition,'' \textit{Proceedings of the IEEE}, vol.~77, no.~2, pp.~257--286, 1989.

\bibitem{b34}
F.-A. Fortin, F.-M. De~Rainville, M.-A. Gardner, M.~Parizeau, and C.~Gagn{\'e}, ``DEAP: Evolutionary algorithms made easy,'' \textit{Journal of Machine Learning Research}, vol.~13, pp.~2171--2175, 2012.

\bibitem{b35}
J.~Wang, H.~Chen, and L.~Zhang, ``AlphaForge: An evolutionary framework for high-capacity formulaic alpha mining,'' \textit{IEEE Transactions on Evolutionary Computation}, vol.~28, no.~4, pp.~1012--1025, 2024.

\bibitem{b36}
Y.~Liu, S.~Wang, and X.~Zhou, ``AlphaQCM: Quantum-inspired genetic programming for cross-sectional alpha discovery,'' \textit{Quantitative Finance}, vol.~25, no.~2, pp.~145--162, 2025.

\end{thebibliography}

\EOD

\end{document}
